% =============================================================================
%  Part VII -- Astrophysics and the strong field
% =============================================================================

Astrophysics probes gravity in two directions away from the laboratory: towards strong fields
(compact objects, black holes, gravitational waves) and towards weak fields over long times and large
distances (galaxies, cosmology). This part states what Causal Physics says in each regime and where
its own content lies.

% =============================================================================
\section{The plateau is general relativity}\label{sec:astro-plateau}
% =============================================================================

On the plateau the order couples to matter as linearised general relativity
(Part~\ref{part:geometric}), and beyond the linear order the coupling to the total conserved
energy--momentum fixes the completion: general relativity. Causal Physics therefore inherits the
astrophysics of general relativity wherever geometry is a reliable description of the order:
\begin{itemize}
\item compact objects and their maximum masses, determined by general relativity and the equation of
  state of matter;
\item black holes with horizons;
\item gravitational waves with two transverse polarisations, propagating at $c$ (to $10^{-15}$ by
  GW170817~\cite{Abbott2017GW}), emitted with the quadrupole formula;
\item the perihelion advance, light deflection, Shapiro delay and frame dragging with their general
  relativistic values.
\end{itemize}
These are consistency requirements, not predictions: a measured deviation from general relativity in any
of them, at scales far above $\ell_k$, would falsify the theory (Sec.~\ref{sec:falsification}).

% =============================================================================
\section{Horizons in the causal picture}\label{sec:astro-horizons}
% =============================================================================

A black hole horizon is, in the causal picture, the limit of reconstructability
(Sec.~\ref{sec:onto-horizons}): the boundary beyond which an outside observer cannot keep an
order-preserving description of the inside. The order itself continues: causal updates go on inside,
and causal continuity is intact. \stat{interpretation}

The entropy of a horizon is a count of the order: the number of horizon molecules is proportional to its
proper area (K.14, Sec.~\ref{sec:reading}). With an entropy $s_0$ per molecule this gives the Bekenstein--Hawking law~\cite{Bekenstein1973,Hawking1975} and ties the discreteness length to Newton's
constant, $\ell\simeq1.075\sqrt{s_0}\,\ell_P$ (Sec.~\ref{sec:bridge}). \stat{tested: K.14; bridge
conditional}

What happens to the order near and inside a horizon, where the load is extreme and the geometric
reading fails, is not described by the present work. Discrete corrections to horizon physics are
expected at the scales $\ell$ and $\ell_k$. \stat{open}

% =============================================================================
\section{Discreteness and non-locality in astrophysical fields}\label{sec:astro-nonlocal}
% =============================================================================

The discrete order leaves two kinds of imprint on long-lived gravitational fields.

\paragraph{Short scales.} The static response of the order deviates from the Einstein--Hilbert form at
wavelengths approaching $\ell_k$: its spin-0 part weakens as
$1-K_{\rm loc}\simeq(k\ell_k)^2[4.0\ln(1/k\ell_k)+0.9]$ and changes sign at a wavelength of
$12.7\,\ell_k$ (Sec.~\ref{sec:bdaction}). For any $\ell_k$ far below astrophysical scales this is
invisible in astrophysics; it matters for laboratory tests of gravity at short distances
(Part~\ref{part:predictions}).

\paragraph{Long times.} The non-local age term weakens the static response of a field that has existed for
a time $T$ by the relative amount $0.1\,k^3\ell_k^2T$ (in units $c=1$). Astrophysical fields are old, and
the effect grows with their age. For the field of the Earth ($L\simeq6400$\,km, $T\simeq4.5$\,Gyr) a
relative deviation below $10^{-2}$ requires $\ell_k\lesssim1$\,mm; for the Sun at $1$\,AU and a Planckian
$\ell_k$ the effect is of order $10^{-79}$. What $T$ is in a real situation (the age of the field, the
age of the source, or the cosmological past of the region) is an open question of the theory; until it
is settled the age term constrains $\ell_k$ only conditionally. The strongest conditional bounds come not from
astrophysics but from short-range tests of gravity in the laboratory (Sec.~\ref{sec:window}). \stat{derived; interpretation of $T$ open}

\paragraph{Gravitational waves.} The same non-locality modifies the propagation of waves with wavelengths
near $\ell_k$. For astrophysical waves ($\lambda\gtrsim10^3$\,km) and any sub-millimetre $\ell_k$ the effect
is negligible. Its form for propagating waves has not been computed. \stat{open}

% =============================================================================
\section{Galaxies and cosmology}\label{sec:astro-cosmo}
% =============================================================================

\paragraph{Galaxies.} Causal Physics in its present form makes no claim about a modified regime of low
accelerations. Whether the discreteness or the non-locality of the order produces one, and at what
scale, is open. The age term weakens the response of long-lived fields at short wavelengths; it does not
strengthen gravity at large distances. In causal set theory the non-locality of the d'Alembertian has been
proposed as an origin of dark matter through off-shell modes~\cite{SaravaniAslanbeigi2015}; whether the
order of Causal Physics, or its flows, have such modes has not been studied. \stat{open}

\paragraph{The cosmological constant.} The number of events in a region fluctuates by about its square
root. Such fluctuations act on the large-scale geometry like a cosmological constant whose magnitude is
set by the four-volume of the observable universe, and for the present universe it is of the order of
the observed value~\cite{Ahmed2004}; in the thermodynamic completion it is an integration constant
(Sec.~\ref{sec:jacobson}). \stat{open}

\paragraph{Expansion.} The growth of the order on the largest scale is the natural candidate for the
expansion of the universe. A growth dynamics that produces it has not been formulated.
\stat{open}
